% Generated draft: requires mathematical and source review.
\documentclass[11pt]{article}
\usepackage[margin=25mm]{geometry}
\usepackage{fontspec}
\setmainfont{Latin Modern Roman}
\usepackage{amsmath,amssymb,mathtools,unicode-math}
\setmathfont{Latin Modern Math}
\usepackage{enumitem,needspace,xurl,hyperref}
\setlength{\parindent}{0pt}
\setlength{\parskip}{5pt}
\setlength{\emergencystretch}{3em}
\newcommand{\N}{\mathbb{N}}
\newcommand{\Z}{\mathbb{Z}}
\newcommand{\Q}{\mathbb{Q}}
\newcommand{\R}{\mathbb{R}}
\newcommand{\C}{\mathbb{C}}
\newcommand{\F}{\mathbb{F}}
\newcommand{\A}{\mathbb{A}}
\newcommand{\PP}{\mathbb{P}}
\newcommand{\E}{\mathbb{E}}
\newcommand{\eps}{\varepsilon}
\newcommand{\dd}{\,\mathrm{d}}
\begin{document}
\section*{UnsolvedMath ID: 800470}
\textbf{Source problem code:} catalogue entry
\subsection{Definitions and mathematical statement}
Let \((N,g)\) be an orientable closed hyperbolic surface, where closed means compact without boundary and the hyperbolic metric has Gaussian curvature \(-1\). Write \(\pi:T^*N\to N\) for the cotangent projection and \(\lambda\) for the canonical one-form, defined by
\[
\lambda_{(q,p)}(v)=p(\dd\pi_{(q,p)}v).
\]
Use the dual metric to define
\[
D_r^*N=\{(q,p)\in T^*N:|p|_g<r\},\qquad D^*N=D_1^*N,
\]
and identify \(N\) with the zero section. The symplectic form in the question is \(d\lambda\).

For precision, use the usual compactly supported convention for the ordinary Hofer--Zehnder capacity. For a symplectic manifold \((U,\omega)\), an admissible Hamiltonian is a function \(H\in C_c^\infty(U,\R_{\geq0})\) that vanishes on a nonempty open set, attains its maximum on a nonempty open set, and whose Hamiltonian flow has no nonconstant periodic orbit of period at most \(1\). The restriction concerns all free homotopy classes. Set
\[
c_{\mathrm{HZ}}(U,\omega)=\sup\{\max H:H\text{ is admissible}\}.
\]
The source does not give its complete admissibility or boundary convention; the open disk convention and period normalization above make the question explicit and require confirmation against the original definition.

\textbf{Question.} Is it true that, for every orientable closed hyperbolic surface \((N,g)\),
\[
c_{\mathrm{HZ}}(D^*N,d\lambda)<\infty?
\]
The registry labels this question open. This is a supplied status label, not an independently verified literature assessment.

\textbf{Source qualifications and context.} The following material records the supplied discussion and its qualifications; it is not a solution or an independent verification of the cited results. The source explicitly describes its two indications toward finiteness and two indications toward infiniteness as not mathematically worked out, and credits discussions with Gabriele Benedetti, Dylan Cant, Stefan Nemirovski, and Alexandru Oancea.

The orientability hypothesis is intentional. The source reports that a non-orientable surface with \(4\mid\chi(N)\) admits a Lagrangian embedding into \((\C^2,\omega_0)\), citing [Giv86]. It then invokes Weinstein's tubular neighborhood theorem to obtain a symplectic embedding of \((D_\varepsilon^*N,d\lambda)\) for some \(\varepsilon>0\), and consequently finiteness of the capacity in that case.

For the relative capacity, impose the additional requirement \(H|_N\equiv\max H\). The source attributes to [Wbr06] the equality
\[
c_{\mathrm{HZ}}(D^*N,N,d\lambda)=\operatorname{sys}(N,g)<\infty,
\]
where \(\operatorname{sys}(N,g)\) is the length of a shortest closed non-contractible geodesic. This relative capacity is distinct from the capacity asked about above.

The source also recalls that displacement energy bounds the ordinary Hofer--Zehnder capacity. It states that displacement of the zero section would imply displacement of a sufficiently small cotangent disk neighborhood and hence finiteness, but that the full zero section of a closed manifold is never Hamiltonian displaceable in its cotangent bundle. In contrast, it asserts that \(N\setminus D_\varepsilon\) is displaceable after an arbitrarily small disk \(D_\varepsilon\subset N\) is removed, and records
\[
c_{\mathrm{HZ}}(D^*(N\setminus D_\varepsilon),d\lambda)<\infty
\]
for every such disk parameter \(\varepsilon>0\). Its proposed argument chooses a Morse function \(f:N\to\R\) with all critical points in \(D_{\varepsilon/2}\), uses \(|\nabla f|\geq\delta>0\) on \(N\setminus D_\varepsilon\), and invokes the flow of \(H(q,p)=\tfrac12|p|_g^2+f(q)\). Boundary conventions for the punctured disk cotangent bundle and the precise displacement construction are not supplied.

For the first indication toward infiniteness, the source describes pseudoholomorphic curves and a form of uniruledness as underlying most known capacity bounds, citing [HV92]. It says that, for closed surfaces whose finiteness is known in its account, this is ensured by a Lagrangian embedding into a ruled complex surface, listing
\[
S^2\subset\mathbb{CP}^1\times\mathbb{CP}^1,\qquad
\mathbb{RP}^2\subset\mathbb{CP}^2,\qquad
T^2\subset\mathbb{CP}^2.
\]
It argues that an orientable surface of genus \(h\geq2\) cannot admit the proposed Lagrangian embedding into a ruled K\"ahler surface \(\Sigma\): such an embedding would give
\[
[N]^2=-\chi(N)=2h-2>0,\qquad
[\omega]\cdot\operatorname{PD}[N]=0,
\]
where \([\omega]\) is the K\"ahler class and \(\operatorname{PD}[N]\) is the Poincar\'e dual of \([N]\). The source invokes the single positive eigenvalue of the intersection form, citing [Bea96, Proposition III.18]. Its use of the term ``ruled'' alongside the listed projective-plane examples requires terminological clarification.

For the second indication toward infiniteness, the source lists finiteness criteria involving a dilation \(\Delta x=1\) [BC25], a nontrivial pair-of-pants product \(x\star y=1\) [Iri14], and vanishing symplectic homology with differential graded local coefficients [AFO17, BDHO-cotangent]. It states that all except the pair-of-pants criterion bound the \(\pi_1\)-sensitive capacity \(c_{\mathrm{HZ}}^0\), defined by excluding fast nonconstant contractible periodic orbits only. It reports
\[
c_{\mathrm{HZ}}^0(D^*N,d\lambda)=\infty
\]
for a hyperbolic surface, using the absence of nonconstant contractible periodic orbits in the geodesic flow of \(E(q,p)=\tfrac12|p|_g^2\). This does not identify the ordinary capacity with its contractible-orbit variant.

The source further invokes a ring isomorphism between cotangent-bundle symplectic homology and free-loop-space homology [Abo15]. With \(\mathcal L N\) the free loop space, \(\pi=\pi_1(N)\), \(\operatorname{Conj}(\pi)\) its set of conjugacy classes, and \(C_\pi(a)\) the centralizer of \(a\in\pi\), it gives, using the \(K(\pi,1)\) property [LO15, Chapter 1.4],
\[
H_*(\mathcal L N;\Z)\cong
\bigoplus_{[a]\in\operatorname{Conj}(\pi)}H_*(K(C_\pi(a),1);\Z)
\cong H_*(N;\Z)\oplus\bigoplus_{[a]\ne[e]}H_*(S^1;\Z),
\]
using \(C_\pi(e)=\pi\) and \(C_\pi(a)\cong\Z\) for \(a\ne e\). Its conclusion that no nontrivial product \(x\star y=1\) exists relies on the degree \(-2\) of the loop product. The meaning of ``nontrivial,'' grading conventions, and the source's assertion that the unit is the only degree \(2\) element require clarification; over \(\Z\), a unit also has integer multiples. The displayed product target in the source also needs correction or explanation.

Finally, the source proposes lower bounds from hyperbolic billiards. It considers an open simply connected \(\Omega\subset N\), describes a lift to \(\mathbb H^2\) with compact closure, and writes
\[
\Omega\text{-}\operatorname{sys}\leq c_{\mathrm{HZ}}(D^*N,d\lambda),
\]
where \(\Omega\text{-}\operatorname{sys}\) is described as the length of a shortest billiard trajectory in \(\Omega\). Compactness of the closure of a chosen lift must be retained as an additional qualification, rather than inferred here from simple connectedness. The source allows a relaxed billiard notion: it only requires closed trajectories arising as limits of approximate solutions in the schemes of [AM11, Voc21] to have length at least \(\Omega\text{-}\operatorname{sys}\). Boundary regularity, existence of a shortest trajectory, and the precise approximation class are unspecified, so this proposed strategy is not a fully formulated auxiliary assertion.
\subsection{Short English statement}
For every orientable compact hyperbolic surface without boundary, ask whether its unit cotangent disk bundle has finite ordinary Hofer--Zehnder capacity. All nonconstant periodic Hamiltonian orbits count, including non-contractible ones. The supplied relative-capacity result and the asserted infinite contractible-orbit capacity concern different quantities. The registry labels the main question open; the surrounding heuristic arguments and billiard proposal retain unresolved qualifications and require mathematical review.
\subsection{Sources}
\begin{itemize}
\item Local UnsolvedMath registry, permanent ID \(800470\), upstream ID \(800470\), exact source problem code \texttt{catalogue entry}, in the supplied AMR Open Problem Lists record. Supplied title: ``Billiards and Quantitative Symplectic Geometry --- Johanna Bimmermann --- Problem 1.'' The registry lists Albach, Barraud et al. as proposers, year 2026, and status open.
\item Supplied primary source: Albach, Barraud et al., \emph{Open Problems in Billiards and Quantitative Symplectic Geometry} (2026), Johanna Bimmermann, Problem 1, \url{https://arxiv.org/abs/2602.12896}. The record states extraction by \texttt{arxiv\_latex} and access on 2026-10-16. This access date and the accompanying literature discussion are supplied provenance, not evidence of independent source access or verification in this draft.
\item Contributor attribution: Johanna Bimmermann. The source credits discussions with Gabriele Benedetti, Dylan Cant, Stefan Nemirovski, and Alexandru Oancea, and explicitly calls the finiteness and infiniteness indications not mathematically worked out.
\item [Giv86]: supplied reference label for the non-orientable Lagrangian embedding assertion. Full bibliographic details were not supplied. Weinstein's tubular neighborhood theorem is also invoked without a bibliographic entry.
\item [Wbr06]: supplied reference label for the relative Hofer--Zehnder capacity and systole equality. Full bibliographic details were not supplied.
\item [HV92]: supplied attribution to Hofer--Viterbo for the pseudoholomorphic-curve approach to capacity bounds. Full bibliographic details were not supplied.
\item [Bea96, Proposition III.18]: supplied reference for the intersection-form assertion. Full bibliographic details were not supplied.
\item [BC25], [Iri14], [AFO17], and [BDHO-cotangent]: supplied reference labels for dilation, pair-of-pants, and local-coefficient symplectic-homology criteria, respectively. The source names Irie in connection with [Iri14]; complete bibliography entries were not supplied.
\item [Abo15]: supplied reference label for the ring isomorphism involving symplectic homology and free-loop-space homology. Full bibliographic details were not supplied.
\item [LO15, Chapter 1.4]: supplied reference for the loop-space discussion of a \(K(\pi,1)\). Full bibliographic details were not supplied.
\item [AM11] and [Voc21]: supplied reference labels for the approximation schemes used in the qualified billiard proposal. Full bibliographic details were not supplied.
\end{itemize}
\end{document}
